\naslov{Računske napake}

Kadar z numerično metodo nekaj izračunamo, ne dobimo točne vrednosti, vendar nek
približek.
\pojem{Absolutno napako} definiramo kot razliko med približkom in točno
vrednostjo:
\[
  d_a = \hat{x} - x.
\]
Po drugi strani je \pojem{relativna napaka} kvocient
\[
  d_r = \frac{\hat{x} - x}{x}.
\]
Približek lahko izrazimo kot $\hat{x} = x (1 + d_r)$.

\vprasanje{Definiraj absolutno in relativno napako.}

Števila predstavljamo s plavajočo vejico, ki je pravzaprav eksponentni zapis
\[
  x = \pm m \cdot b^e,
\]
kjer je $m$ \pojem{mantisa}, zapisana kot $m = 0.c_1c_2\ldots c_t$ za $c_i \in
\{0, \ldots, b-1\}$, število $b$ je \pojem{baza} zapisa, $e$ pa
\pojem{eksponent} v mejah $L \le e \le U$.
Števila običajno zapišemo \pojem{normalizirana}, torej s $c_1 \ne 0$.
V primeru najnižje možne potence dovoljujemo tudi \pojem{subnormalizirana}
števila, kjer je $c_1 = 0$.
Predstavljiva števila v takšnem zapisu označujemo s $P(b, t, L, U)$.

V standardu IEEE imamo dve števili:
\begin{itemize}
\item Enojni zapis: $P(2, 24, -125, 128)$,
\item Dvojni zapis: $P(2, 53, -1021, 1023)$.
\end{itemize}

\vprasanje{Kaj je $P(b, t, L, U)$? Kakšne vrednosti imata \texttt{float} in
  \texttt{double}?}

Pri zaokoževanju številu odrežemo decimalke za neko vrednostjo, in po potrebi
prištejemo $b^{-t}$.
Boljšega od teh približkov označimo s $\fl(x)$.

\begin{izrek}
  Če za $x$ velja, da $\abs{x}$ leži na intervalu med najmanjšim in največjim
  pozitivnim predstavljivim normaliziranim številom, potem velja
  \[
    \frac{\abs{\fl(x) - x}}{\abs{x}} \le u,
  \]
  za \pojem{osnovno zaokrožitveno napako} $u = \pol b^{1-t}$.
\end{izrek}

\vprasanje{Kaj je osnovna zaokrožitvena napaka? Povej izrek.}

Standard IEEE zagotovlja omejeno napako tudi pri osnovnih operacijah:
\begin{itemize}
\item $\fl(x \oplus y) = (x \oplus y) (1 + \delta)$ za $\abs{\delta} \le u$ za
  osnovne operacije $+, -, \cdot, /$,
\item $\fl(\sqrt{x}) = \sqrt{x} (1 + \delta)$ za $\abs{\delta} \le u$.
\end{itemize}

Drug vir napak je občutljivost problema, ki ni povezana z numeriko.
Obravanavamo vprašanje, kako se pri majhni spremembi v vhodnih podatkih spremeni
pravilni odgovor.
Za zvezno odvedljivo $f$ lahko absolutno občutljivost merimo z odvodom
\[
  \abs{f(x+\delta x) - f(x)} \approx \abs{f'(x)} \abs{\delta x}.
\]

Poznamo tri vrste napak, ki skupaj sestavljajo celotno napako:
\begin{itemize}
\item Neodstranljiva napaka: napaka zaradi zaokroževanja podatkov
\item Napaka metode: nenatančnost metode
\item Zaokrožitvena napaka: napaka zaradi zaokroževanja znotraj metode
\end{itemize}

\vprasanje{Katere vrste napak poznamo?}